% Synthetic educational data only. Replace every field for your own thesis.
\hitsetup{
  statesecrets={公开},
  natclassifiedindex={TP391},
  intclassifiedindex={004},
  ctitlecover={学位论文模板使用示例},
  ctitle={学位论文模板使用示例},
  cxueke={工学},
  csubject={计算机科学与技术},
  caffil={示例学院},
  cauthor={示例学生},
  csupervisor={示例导师教授},
  cassosupervisor={},
  ccosupervisor={},
  cstudentid={DEMO-ONLY},
  cdate={2026 年 10 月},
  firstpagecdate={2026 年 10 月},
  etitle={An Example of Using the Thesis Template},
  exueke={Engineering},
  esubject={Computer Science and Technology},
  eaffil={Example School},
  eauthor={Example Student},
  esupervisor={Professor Example},
  eassosupervisor={},
  ecosupervisor={},
  edate={October, 2026},
  ckeywords={学位论文,排版,模板},
  ekeywords={thesis,typesetting,template},
}
\begin{cabstract}
本示例用于说明模板的基本结构和独立选项，不代表实际研究成果。请按学位论文撰写规范改写题目、摘要、正文与后置材料，逐项核对本人的学位类型、学院要求和提交渠道。
\end{cabstract}
\begin{eabstract}
This synthetic example illustrates the document structure and independent
template options. It does not report research results. Replace all metadata
and content, then check the requirements for your own degree and submission.
\end{eabstract}
